\documentclass[12pt,a4paper,bibliography=totocnumbered,listof=totocnumbered]{scrartcl}
\usepackage[ngerman]{babel}
\usepackage[utf8]{inputenc}
\usepackage{amsmath}
\usepackage{amsfonts}
\usepackage{amssymb}
\usepackage{graphicx}
\usepackage{fancyhdr}
\usepackage{tabularx}
\usepackage{geometry}
\usepackage{setspace}
\usepackage[right]{eurosym}
\usepackage[printonlyused]{acronym}
\usepackage{subfig}
\usepackage{floatflt}
\usepackage[usenames,dvipsnames]{color}
\usepackage{colortbl}
\usepackage{paralist}
\usepackage{array}
\usepackage{titlesec}
\usepackage{parskip}
\usepackage[right]{eurosym}
%\usepackage{url}
%\usepackage{wrapfig}
%\usepackage[backend=biber, citestyle=authoryear]{biblatex}

\usepackage[subfigure,titles]{tocloft}
\usepackage[pdfpagelabels=true]{hyperref}

\usepackage{listings}
\lstset{basicstyle=\footnotesize, captionpos=b, breaklines=true, showstringspaces=false, tabsize=2, frame=lines, numbers=left, numberstyle=\tiny, xleftmargin=2em, framexleftmargin=2em}
\makeatletter
\def\l@lstlisting#1#2{\@dottedtocline{1}{0em}{1em}{\hspace{1,5em} Lst. #1}{#2}}
\makeatother

\geometry{a4paper, top=27mm, left=30mm, right=20mm, bottom=35mm, headsep=10mm, footskip=12mm}

\hypersetup{unicode=false, pdftoolbar=true, pdfmenubar=true, pdffitwindow=false, pdfstartview={FitH},
	pdftitle={Projektkurs Dokumentation},
	pdfauthor={Jakob Fleischer, Robin Schlaak, Lars Schmalbach},
	pdfsubject={Dokumentation},
	pdfcreator={\LaTeX\ with package \flqq hyperref\frqq},
	pdfproducer={pdfTeX \the\pdftexversion.\pdftexrevision},
	pdfkeywords={Dokumentation},
	pdfnewwindow=true,
	colorlinks=true,linkcolor=black,citecolor=black,filecolor=magenta,urlcolor=black}
\pdfinfo{/CreationDate \today}

%\bibliography{bibo.bib}
%\addbibresource{bibo.bib}

\begin{document}

\titlespacing{\section}{0pt}{12pt plus 4pt minus 2pt}{-6pt plus 2pt minus 2pt}

% Kopf- und Fusszeile
\renewcommand{\sectionmark}[1]{\markright{#1}}
\renewcommand{\leftmark}{\rightmark}
\pagestyle{fancy}
\lhead{}
\chead{}
\rhead{\thesection\space\contentsname}
\lfoot{Dokumentation unseres Projekts\newline im Projektkurs Mathe-Physik-Informatik}
\cfoot{}
\rfoot{\ \linebreak Seite \thepage}
\renewcommand{\headrulewidth}{0.4pt}
\renewcommand{\footrulewidth}{0.4pt}

% Vorspann
\renewcommand{\thesection}{\Roman{section}}
\renewcommand{\theHsection}{\Roman{section}}
\pagenumbering{Roman}

% ----------------------------------------------------------------------------------------------------------
% Titelseite
% ----------------------------------------------------------------------------------------------------------
\thispagestyle{empty}
\begin{center}
	\includegraphics[scale=2]{Bilder/aeg-logo.png}\\
	
	\Large
	\textbf{Albert-Einstein-Gymnasium}\\
	\vspace*{2cm}
	\vspace*{2cm}
	\Huge
	\textbf{Dokumentation}\\
	\vspace*{0.5cm}
	\large
	zu unserem Projekt im Projektkurs\\
	\vspace*{1cm}
	\Huge
	\textbf{Mathe-Physik-Informatik}\\
	\vspace*{2cm}
	
	\vfill
	\normalsize
	\newcolumntype{x}[1]{>{\raggedleft\arraybackslash\hspace{0pt}}p{#1}}
	\begin{tabular}{x{6cm}p{7.5cm}}
		\rule{0mm}{5ex}\textbf{Autoren:} & Jakob Fleischer\newline Robin Schlaak \newline
		Lars Schmalbach \\ 
		\rule{0mm}{5ex}\textbf{Prüfer:} & Herr Thomas Bachran \\ 
		\rule{0mm}{5ex}\textbf{Abgabedatum:} & 21.06.2017 \\ 
	\end{tabular} 
\end{center}
\pagebreak

% ----------------------------------------------------------------------------------------------------------
% Abstract
% ----------------------------------------------------------------------------------------------------------
\setcounter{page}{1}
\onehalfspacing
\titlespacing{\section}{0pt}{12pt plus 4pt minus 2pt}{2pt plus 2pt minus 2pt}
\rhead{KURZFASSUNG}
\section{Umschreibung des Projekts}
\emph{Jakob}
...

\vspace{-1,2em}
\titlespacing{\section}{0pt}{12pt plus 4pt minus 2pt}{-6pt plus 2pt minus 2pt}

\pagebreak

% ----------------------------------------------------------------------------------------------------------
% Verzeichnisse
% ----------------------------------------------------------------------------------------------------------
% TODO Typ vor Nummer
\renewcommand{\cfttabpresnum}{Tab. }
\renewcommand{\cftfigpresnum}{Abb. }
\settowidth{\cfttabnumwidth}{Abb. 10\quad}
\settowidth{\cftfignumwidth}{Abb. 10\quad}

\titlespacing{\section}{0pt}{12pt plus 4pt minus 2pt}{2pt plus 2pt minus 2pt}
\singlespacing
\rhead{INHALTSVERZEICHNIS}
\renewcommand{\contentsname}{II Inhaltsverzeichnis}
\phantomsection
\addcontentsline{toc}{section}{\texorpdfstring{II \hspace{0.35em}Inhaltsverzeichnis}{Inhaltsverzeichnis}}
\addtocounter{section}{1}
\tableofcontents
\pagebreak


\rhead{VERZEICHNISSE}
\listoffigures
\pagebreak
\listoftables
\pagebreak
\renewcommand{\lstlistlistingname}{Listing-Verzeichnis}
{\labelsep2cm\lstlistoflistings}
\pagebreak

% ----------------------------------------------------------------------------------------------------------
% Abkürzungen
% ----------------------------------------------------------------------------------------------------------
%\section{Abkürzungsverzeichnis}
%\begin{acronym}[OSGi] % längste Abkürzung steht in eckigen Klammern
%	\setlength{\itemsep}{-\parsep} % geringerer Zeilenabstand
%	\acro{OSGi}{Open Service Gateway initiative}
%\end{acronym}
%\newpage

% ----------------------------------------------------------------------------------------------------------
% Inhalt
% ----------------------------------------------------------------------------------------------------------
% Abstände Überschrift
\titlespacing{\section}{0pt}{12pt plus 4pt minus 2pt}{-6pt plus 2pt minus 2pt}
\titlespacing{\subsection}{0pt}{12pt plus 4pt minus 2pt}{-6pt plus 2pt minus 2pt}
\titlespacing{\subsubsection}{0pt}{12pt plus 4pt minus 2pt}{-6pt plus 2pt minus 2pt}

% Kopfzeile
\renewcommand{\sectionmark}[1]{\markright{#1}}
\renewcommand{\subsectionmark}[1]{}
\renewcommand{\subsubsectionmark}[1]{}
\lhead{Kapitel \thesection}
\rhead{\rightmark}

\onehalfspacing
\renewcommand{\thesection}{\arabic{section}}
\renewcommand{\theHsection}{\arabic{section}}
\setcounter{section}{0}
\pagenumbering{arabic}
\setcounter{page}{1}

% ----------------------------------------------------------------------------------------------------------
% Einleitung
% ----------------------------------------------------------------------------------------------------------
\section{Kapitel: Herangehensweise}
\emph{Jakob}\\
Bei der Planung und der Herangehensweise an unser Projekt sind wir vermehrt auf Probleme gestoßen, welche in diesem Kapitel dargestellt werden.
...

\subsection{ursprüngliche Vorstellung und Konsequenzen}
\emph{Lars}\\
Unsere Unerfahrenheit in der Projektarbeit hat uns zu einem überstürzten und unüberlegten Projekteinstieg getrieben, bei dem wir uns in Themen eingearbeitet haben, ohne uns einen genauen Plan zu erstellen.
...\\
Zu spät haben wir gemerkt, dass unsere bisherige Arbeit weder zielführend, noch in irgendeiner Weise sinnvoll war, weshalb wir uns zusammengesetzt und intensiv beraten haben.
....

\subsection{verspätete Planung}
\emph{Jakob}\\
...

\subsection{Zusammenarbeit und Kommunikation}
\emph{Jakob}\\
....

\pagebreak
% ----------------------------------------------------------------------------------------------------------
% Grundlagen
% ----------------------------------------------------------------------------------------------------------

\section{Kapitel: Grundlagen}
\emph{Jakob \emph{\&} Lars}\\
Im zweiten Kapitel dieser Dokumentation wird genauer auf die Themen eingegangen, in welche wir uns im Laufe des Projektkurses hineingearbeitet haben. Es wird nicht nur auf für das Projekt relevante Techniken, sondern auch auf solche, in die wir uns zwar eingearbeitet haben, sie jedoch nicht verwendet haben, eingegangen.

\subsection{\LaTeX{}}
\emph{Robin}\\
...

\subsection{XAMPP}
\emph{Robin}\\
...

\subsection{HTML}
\label{sec:HTML}
\emph{Lars}\\
HTML ist die Abkürzung für "Hypertext Markup Language", was zu deutsch "Hypertext-Auszeichnungssprache" heißt. Es ist eine Programmiersprache, mit der man den Aufbau von Internetseiten bestimmt. Solche HTML-Dokumente stellen die Grundlage für das World Wide Web dar und werden von Browsern dargestellt.\cite{HTML}\cite{Hypertext_Markup_Language} Sie bestehen in der Regel aus drei Teilen.\cite{HTML/Dokumentstruktur_und_Aufbau} Der erste Teil ist die Dokumenttyp-Deklaration. In ihr werden Angaben zur verwendeten HTML-Version gegeben. Im "head", welcher den zweiten Teil darstellt, werden unsichtbare Kopfdaten, zum Beispiel der Titel der Seite oder andere, für den menschlichen Betrachter der Webseite zunächst nicht sichtbare, Informationen zur korrekten Darstellung des sichtbaren Teils der Webseite, angegeben.\cite{HTML/Kopfdaten} Anzuzeigende Inhalte werden in den "body", den dritten Teil, geschrieben. Hier werden also sämtliche Texte, Verweise, Grafiken und so weiter eingefügt, die auf der Webseite sichtbar sein sollen.\cite{HTML/Dokumentstruktur_und_Aufbau}

\subsubsection{Frames}
Eine hilfreiche Technik, die auch bei uns ihren Einsatz gefunden hat, heißt Frames. Diese Technik wurde 1996 von Netscape eingeführt. Mit ihr kann man mehrere Dateien gleichzeitig auf dem Bildschirm anzeigen lassen.\cite{HTML/Frames} Sie wurde jedoch im Oktober 2014 mit HTML5\cite{HTML5} aus dem Standard entfernt, da sie entscheidende Nachteile aufweist. Aufgrund des Verwendungszweckes unserer Webseite benutzen wir Frames, obwohl empfohlen wird, Server-seitig andere Techniken zum Auslagern von Teilen der Seite zu benutzen.\cite{HTML/Frames} Die stärksten Argumente waren, die simple Handhabung und die guten Gestaltungsmöglichkeiten mit dieser Technik.

\subsubsection{Formulare}
\label{sec:Formulare}
Formulare sind ein Element von HTML [\ref{sec:HTML}], das es ermöglicht Daten zu erfassen und über das Hypertext Transfer Protocol per XMLHttpRequest, HTTP-GET oder HTTP-POST zur Verarbeitung an einen Server zu senden.\cite{Webformular} In unserem Programm kam letzteres zum Einsatz. Man kann in HTML zwar Formulare definieren und erstellen, für eine Verarbeitung und Auswertung der Eingaben ist jedoch eine andere Programmiersprache, wie zum Beispiel Javascript [\ref{sec:JavaScript}] oder PHP [\ref{sec:PHP}] nötig.\cite{HTML/Formulare/Form} In unserem Fall haben wir Formulare in Form von Anmeldeformularen, Suchfiltern und auch als Möglichkeit zum Hochladen von Dateien eingesetzt.

\subsubsection{JavaScript}
\label{sec:JavaScript}
Bei JavaScript handelt es sich um eine interpretierende Programmier- beziehungsweise Skriptsprache, die 1995 von Netscape entwickelt wurde.\cite{JavaScript}\cite{wiki/JavaScript} JavaScript ist sehr verbreitet, da sich in allen modernen Browsern Interpreter für die Sprache befinden. Es wird hauptsächlich Client-seitig verwendet und ermöglicht es, dynamischen Einfluss auf Webseiten zu nehmen.\cite{JavaScript-JS} Für unser Programm kam die Sprache besonders häufig aufgrund des, durch sie ermöglichten, einfachen Umgangs mit Variablen und Funktionen zum Einsatz.

\subsection{PHP}
\label{sec:PHP}
\emph{Robin}\\
...

\subsection{Apache}
\emph{Lars}\\
...

\pagebreak
\subsection{Editoren und integrierte Entwicklungsumgebungen}
\emph{Lars}\\
Editoren werden zum Schreiben von Texten, wie zum Quellcodes, benutzt. gute Editoren helfen das Programmieren zu vereinfachen, indem sie gewisse Schlüsselwörter, sowie Befehle, farblich hervorheben, eine Autovervollständigung, eine Such- und Ersetzfunktion und indem sie den Quellcode automatisch einrücken, sowie eine Schnittstelle für Plugins darstellen.\cite{Texteditor} \\
Neben Editoren gibt es auch integrierte Entwicklungsumgebungen. Diese bestehen aus einer Sammlung an Computerprogrammen, mit denen es möglich ist Software ohne die Verwendung vieler einzelner Programme zu entwickeln. Durch sie werden nicht nur Tippfehler verhindert, sondern auch Arbeitsschritte und somit Zewit bei der Softwareentwicklung geplant.\cite{Integrierte_Entwicklungsumgebung}\cite{Medienbruch}
Bei der Entwicklung unseres Programms kamen die Editoren und integrierten Entwicklungsumgebungen Eclipse [\ref{sec:Eclipse}] (genau genommen Eclipse Neon IDE), Notepad++ [\ref{sec:Notepad++}] und Texmaker [\ref{sec:Texmaker}] zum Einsatz.

\subsubsection{Eclipse}
\label{sec:Eclipse}
\emph{Robin}\\
...

\subsubsection{Notepad ++}
\label{sec:Notepad++}
\emph{Robin}\\
...

\subsubsection{Texmaker}
\label{sec:Texmaker}
\emph{Robin}\\
...

\subsection{GitHub}
\emph{Jakob}\\
Durch die Arbeit mit GitHub, wird das gemeinsame, nicht zwingend parallele, Arbeiten möglich. GitHub verwaltet die Quellcodes, welche die Benutzer hochladen so, dass jeder, der über Zugriff auf das Projekt verfügt, dieses weiter führen kann. -Lars

\subsection{Dropzone}
\emph{Lars}\\
Ursprünglich beschreibt der Begriff \glqq Dropzone\grqq{} einen geheimen Speicherort für maschinell gestohlene Daten, wie zum Beispiel Passwörter und Kontodaten.\cite{wiki/Dropzone} \\
In unserem Fall jedoch ist die Dropzone ein Script, welches in JavaScript [\ref{sec:JavaScript}] geschrieben wurde. Dieses dient zur einfachen Gestaltung des Datei-Uploads mittels HTML. Es soll die optische Anpassung des Eingabefelds vereinfachen. 
Aufgrund ihrer Komplexität, welche sich leider erst nach und nach herauskristallisierte, haben wir uns schlussendlich gegen die Dropzone und für einen herkömmlichen Datei-Upload mittels HTML-Formular [\ref{sec:Formulare}] entschieden.

\subsection{Java-Servlet}
\emph{Jakob}\\
Im Folgenden werde ich erläutern, wie man ein Servlet in Eclipse über einen Apache Tomcat Server erstellen kann.

\subsubsection{Einrichtung in Eclipse}
Vorausgesetzt für eine erfolgreiche Einrichtung ist die bereits erfolgte Installation von Eclipse Neo in der Java EE Version und Apache Tomcat v 9.0. 
Nachdem  diese erfolgt ist, ist es nun möglich in Eclipse einen neuen Server einzurichten, wobei man unter dem Ordner Apache die installierte Tomcat Version vorfindet, welche hier  als Servertyp verwendet wird.
Nachdem dies geschehen ist, richten sie ein Dynamic Web Project ein und innerhalb desselben erstellen sie im Java Resources/src  Ordner ein neues Servlet, indem sie nun programmieren können.

Wenn nun aber die Fehlermeldung  „cannot resolve“ im Bezug auf die vorgegebenen import-Zeilen auftaucht müssen sie außerdem noch die servlet-api.jar downloaden und diese dann in den Eigenschaften ihres Dynamic Web Projects unter Java Build Path in der Kategorie Libraries mit der Funktion „Add External JARs“ hinzufügen. 

...

\subsubsection{Vorteile}
...

\subsubsection{Nachteile}
...

\subsection{JSON}
\emph{Jakob}\\
Ein weiteres Thema, in welches ich mich eingearbeitet habe war die Erstellung und das Auslesen von Dateien
des Formates JSON mit Eclipse.
...

\subsubsection{Erstellen}
...

\subsubsection{Auslesen}
...

\subsubsection{Vorteile}
...

\subsubsection{Nachteile}
...

\subsection{MYSQL}
\emph{Jakob}\\
Nachdem obige Thematiken angesichts unseres erneuerten Projektplans irrelevant geworden waren, habe ich mich 
mit MYSQL beschäftigt.
...

\subsubsection{Vorteile}
...

\subsubsection{Datenbank-Entwurf}
Im Folgenden möchte ich auf das Entwerfen einer MYSQL-Datenbank eingehen.
...

\subsubsection{Verwaltung einer Datenbank mit PHP}
Um eine Datenbank in unser Projekt zweckmäßig einzurichten, musste ich mich in eine MYSQL-Einbindung mittels PHP einlesen, welche ich nun darstellen möchte.
...

\pagebreak

% ----------------------------------------------------------------------------------------------------------
% Kapitel 3 - Projekt
% ----------------------------------------------------------------------------------------------------------
\section{Projekt}
\emph{Lars}\\
In diesem Kapitel wird nun unser Projekt dargestellt, indem zunächst das Programm im Gesamten beschrieben wird. Daraufhin wird die Arbeit von Lars am GUI, die Arbeit von Robin am Interface und File-System und zuletzt die Arbeit von Jakob am Interface und an der Datenbank dargestellt. 
...

\subsection{File-System-Management-System}
\emph{Lars}\\
Darstellung Ergebnis...

\subsection{GUI}
\emph{Lars}\\
Darstellung GUI...

\subsection{Engine I: Interface und File-System}
\emph{Robin}\\
Darstellung Interface und File-System...

\subsection{Engine II: Interface und Datenbank}
\emph{Jakob}\\
Darstellung Interface und Datenbank...
\pagebreak
% ----------------------------------------------------------------------------------------------------------
% Kapitel
% ----------------------------------------------------------------------------------------------------------
\section{Fazit}
Im Folgenden wird jeder von uns ein Fazit geben, indem er darstellt, was er aus dem Projekt gelernt hat, was ihm gefallen hat, Sachen die er gut findet und so weiter.

\subsection{}
\emph{Lars}\\
...

\subsection{}
\emph{Robin}\\
...

\subsection{}
\emph{Jakob}\\
...

\pagebreak
% ----------------------------------------------------------------------------------------------------------
% Literatur
% ----------------------------------------------------------------------------------------------------------

%\addcontentsline{tocnumbered}{section}{Literaturverzeichnis}

\begin{thebibliography}{999}

%\bibitem[Anzeigename]{Stichwort, dass man mit \cite aufruft}
%\url{link zur Quelle}, Zugriff: 16.06.2017

\bibitem{HTML}
\url{https://wiki.selfhtml.org/wiki/HTML}, \\
Zugriff: 15.06.2017

\bibitem{Hypertext_Markup_Language}
\url{https://de.wikipedia.org/wiki/Hypertext\_Markup\_Language},\\
Zugriff: 15.06.2017

\bibitem{HTML/Dokumentstruktur_und_Aufbau}
\url{https://wiki.selfhtml.org/wiki/HTML/Dokumentstruktur\_und\_Aufbau}, \\
Zugriff: 15.06.2017

\bibitem{HTML/Kopfdaten}
\url{https://wiki.selfhtml.org/wiki/HTML/Kopfdaten}, \\
Zugriff: 15.06.2017

\bibitem{HTML/Frames}
\url{https://wiki.selfhtml.org/wiki/HTML/Frames}, \\
Zugriff: 15.06.2017

\bibitem{HTML5}
\url{https://de.wikipedia.org/wiki/HTML5}, \\
Zugriff: 15.06.2017

\bibitem{Webformular}
\url{https://de.wikipedia.org/wiki/Webformular}, \\
Zugriff: 15.06.2017

\bibitem{JavaScript}
\url{http://www.searchenterprisesoftware.de/definition/JavaScript}, \\
Zugriff: 15.06.2017

\bibitem{wiki/JavaScript}
\url{https://de.wikipedia.org/wiki/JavaScript}, \\
Zugriff: 15.06.2017

\bibitem{JavaScript-JS}
\url{http://www.itwissen.info/JavaScript-JavaScript-JS.html}, \\
Zugriff: 15.06.2017

\bibitem{HTML/Formulare/Form}
\url{https://wiki.selfhtml.org/wiki/HTML/Formulare/Form}, \\
Zugriff: 15.06.2017

\bibitem{Integrierte_Entwicklungsumgebung}
\url{https://de.wikipedia.org/wiki/Integrierte_Entwicklungsumgebung}, \\
Zugriff: 16.06.2017

\bibitem{Texteditor}
\url{https://de.wikipedia.org/wiki/Texteditor}, \\
Zugriff: 16.06.2017

\bibitem{Medienbruch}
\url{https://de.wikipedia.org/wiki/Medienbruch}, \\
Zugriff: 16.06.2017

\bibitem{wiki/Dropzone}
\url{https://de.wikipedia.org/wiki/Dropzone}, \\
Zugriff: 16.06.2017



%\bibitem{}
%\url{}, \\
%Zugriff: 16.06.2017


\end{thebibliography}
%\section{Quellenverzeichnis}
%\printbibliography
%\renewcommand\refname{Quellenverzeichnis}
%\bibliographystyle{alpha}
%\bibliography{bibo}
\pagebreak


% ----------------------------------------------------------------------------------------------------------
% Anhang
% ----------------------------------------------------------------------------------------------------------
\pagenumbering{Roman}
\setcounter{page}{1}

\lhead{Anhang \thesection}
\begin{appendix}

\section*{Anhang}
\phantomsection
\addcontentsline{toc}{section}{Anhang}
\addtocontents{toc}{\vspace{-0.5em}}


\end{appendix}


\newpage
\thispagestyle{empty}
\begin{center}
	\vspace*{5em}
	\huge\textbf{Erklärung}\\
\end{center}
\vspace{2em}
Hiermit versichere ich, dass ich meine Dokumentation selbständig verfasst und keine anderen als die angegebenen Quellen und Hilfsmittel benutzt habe.

\vspace{4em}
\begin{minipage}{\linewidth}
	\begin{tabular}{p{15em}p{15em}}
		Datum: &  .......................................................\\
		& \centering (Unterschrift)\\
	\end{tabular}
\end{minipage}

\end{document}

%Die Quellen befinden sich in der Datei \textit{bibo.bib}. Ein Buch- und eine Online-Quelle sind beispielhaft eingefügt. [Vgl. \cite{buch}, \cite{online}]
%Abkürzungen lassen sich natürlich auch nutzen (\ac{OSGi}). Weiter oben im Latex-Code findet sich das %Verzeichnis.


%Beispiele
%
%
%----------Listings----------
%
%
% \ref{lst:arduino}.
%
%\vspace{1em}
%\begin{lstlisting}[caption=Arduino Beispielprogramm, label=lst:arduino]
%int ledPin = 13;
%void setup() {
%    pinMode(ledPin, OUTPUT);
%}
%void loop() {
%    digitalWrite(ledPin, HIGH);
%    delay(500);
%    digitalWrite(ledPin, LOW);
%    delay(500);
%}
%\end{lstlisting}
%
%
%
%----------Aufzählung----------
%
%
%\begin{compactitem}
%	\item Nur
%	\item ein
%	\item Beispiel.
%%\end{compactitem}
%
%
%
%----------Tabelle----------
%
%
%\vspace{1em}
%\begin{table}[!h]
%	\centering
%	\begin{tabular}{|l|l|l|}
%		\hline
%		\textbf{Name} & \textbf{Name} & \textbf{Name}\\
%		\hline
%		1 & 2 & 3\\
%		\hline
%		4 & 5 & 6\\
%		\hline
%		7 & 8 & 9\\
%		\hline
%	\end{tabular}
%	\caption{Beispieltabelle}
%	\label{tab:beispiel}
%\end{table}
%
%
%
%----------Grafik/Abbildung----------
%
%
%\vspace{1em}
%\begin{minipage}{\linewidth}
%	\centering
%	\includegraphics[width=0.7\linewidth]{Bilder/layering-osgi.png}
%	\captionof{figure}[OSGi Architektur]{OSGi Architektur\footnotemark }
%	\label{fig:osgi}
%\end{minipage}
%\footnotetext{Quelle: \url{http://www.osgi.org/Technology/WhatIsOSGi}}



%section*{}  -> Section, die nicht im Inhaltsverzeichnis aufgeführt wird



